\documentclass[11pt,a4paper]{article}

% ── Packages ─────────────────────────────────────────────────────────────────
\usepackage[utf8]{inputenc}
\usepackage[T1]{fontenc}
\usepackage[margin=2.5cm]{geometry}
\usepackage{amsmath,amssymb}
\usepackage{booktabs}
\usepackage{array}
\usepackage{longtable}
\usepackage{graphicx}
\usepackage{xcolor}
\usepackage{listings}
\usepackage{hyperref}
\usepackage{enumitem}
\usepackage{fancyhdr}
\usepackage{titlesec}
\usepackage{parskip}
\usepackage{microtype}

% ── Line-breaking tolerance ──────────────────────────────────────────────────
% Many paragraphs below contain long inline \texttt identifiers (snake_case
% API field names) that cannot be hyphenated; relax justification slightly
% rather than letting them overflow the text width.
\sloppy
\emergencystretch=3em

% ── Colors ───────────────────────────────────────────────────────────────────
\definecolor{codebg}{RGB}{245,245,245}
\definecolor{codekw}{RGB}{0,90,150}
\definecolor{linkblue}{RGB}{0,90,160}

% ── Hyperref setup ───────────────────────────────────────────────────────────
\hypersetup{
    colorlinks=true,
    linkcolor=linkblue,
    urlcolor=linkblue,
    citecolor=linkblue,
    pdftitle={MotorAudit Technical Documentation},
    pdfauthor={MotorAudit Project},
    pdfsubject={Energy audit platform for industrial electric motors},
}

% ── Code listing style ───────────────────────────────────────────────────────
\lstdefinestyle{code}{
    backgroundcolor=\color{codebg},
    basicstyle=\ttfamily\small,
    breaklines=true,
    frame=single,
    rulecolor=\color{gray!40},
    columns=fullflexible,
    keepspaces=true,
}
\lstset{style=code}

% ── Headers/footers ──────────────────────────────────────────────────────────
\pagestyle{fancy}
\fancyhf{}
\fancyhead[L]{MotorAudit Documentation}
\fancyhead[R]{\leftmark}
\fancyfoot[C]{\thepage}
\renewcommand{\headrulewidth}{0.4pt}

\titleformat{\section}{\Large\bfseries}{\thesection}{0.6em}{}
\titleformat{\subsection}{\large\bfseries}{\thesubsection}{0.6em}{}

\newcommand{\code}[1]{\texttt{#1}}
\newcommand{\apiendpoint}[2]{\texttt{\textbf{#1}} \texttt{#2}}

% ── Title page data ──────────────────────────────────────────────────────────
\title{\vspace{-2cm}\textbf{MotorAudit}\\[0.3em]
\large Technical Documentation \\[0.2em]
\normalsize Energy Audit Platform for Industrial Electric Motors}
\author{MotorAudit Project}
\date{Version 1.0.0 \\ \today}

\begin{document}

\maketitle
\thispagestyle{empty}

\begin{abstract}
\noindent
MotorAudit is a web-based engineering platform that helps plant engineers and
facility managers quantify the energy-saving potential of three common
industrial interventions: replacing an inefficient motor, adding a variable
frequency drive (VFD) to a throttled pump or fan, and diagnosing power-quality
problems from voltage/current waveforms. The three modules share a common IEC
motor efficiency database and can be consolidated into a single PDF audit
report. This document describes the system architecture, the engineering
formulas implemented in each module, the REST API, and the assumptions and
limitations that apply to the results.
\end{abstract}

\thispagestyle{empty}
\newpage

\tableofcontents
\newpage

% ══════════════════════════════════════════════════════════════════════════
\section{Introduction}
% ══════════════════════════════════════════════════════════════════════════

MotorAudit combines three engineering modules under a single user interface:

\begin{enumerate}[label=\textbf{\arabic*.}]
    \item \textbf{Motor Replacement \& Payback} --- estimates the annual
          energy, cost, and CO\textsubscript{2} savings from replacing an
          existing motor with a higher-efficiency model, and the resulting
          simple payback period.
    \item \textbf{VFD Sizing \& Savings} --- estimates the energy savings
          from installing a variable frequency drive on a throttled pump,
          fan, or conveyor, based on a user-supplied speed/time load profile.
    \item \textbf{Power Quality Analyzer} --- computes total harmonic
          distortion (THD), power factor, and phase unbalance from
          voltage/current waveform samples (uploaded as CSV, submitted as
          JSON, or generated from built-in presets), and flags out-of-limit
          conditions with corrective recommendations.
\end{enumerate}

All three modules can optionally be enriched with natural-language
engineering commentary generated by a large language model (Google Gemini),
and the results of all three can be merged into a single downloadable PDF
audit report.

Results throughout the application are explicitly labelled as
\textbf{engineering estimates} rather than certified measurements; see
Section~\ref{sec:assumptions} for the full list of assumptions and
limitations.

% ══════════════════════════════════════════════════════════════════════════
\section{System Architecture}
% ══════════════════════════════════════════════════════════════════════════

MotorAudit is a two-tier web application: a Python backend that performs all
engineering calculations and a TypeScript/React frontend that collects inputs
and visualizes results.

\subsection{Technology Stack}

\begin{table}[h]
\centering
\begin{tabular}{@{}ll@{}}
\toprule
\textbf{Layer} & \textbf{Technology} \\
\midrule
Backend framework      & FastAPI (Python), served by Uvicorn \\
Data validation         & Pydantic / \texttt{pydantic-settings} \\
Numerical computation   & NumPy \\
PDF generation          & ReportLab \\
AI insights             & Google Gemini (\texttt{gemini-3-flash-preview}) \\
Frontend framework      & Next.js (App Router) / React / TypeScript \\
Styling                 & Tailwind CSS \\
Testing                 & pytest \\
\bottomrule
\end{tabular}
\caption{Core technology stack.}
\end{table}

\subsection{Repository Layout}

\begin{lstlisting}[style=code]
motoraudit/
|-- backend/
|   |-- app/
|   |   |-- main.py              # FastAPI app entry point
|   |   |-- config.py            # Settings (pydantic-settings)
|   |   |-- db/seed.py           # IEC motor efficiency table
|   |   |-- schemas/             # Pydantic request/response models
|   |   |-- core/
|   |   |   |-- motor_db.py      # Efficiency interpolation
|   |   |   |-- replacement_calc.py
|   |   |   |-- vfd_calc.py
|   |   |   |-- dsp.py           # FFT, THD, PF, unbalance
|   |   |   |-- gemini.py        # AI insight generation
|   |   |   `-- report_gen.py    # ReportLab PDF generator
|   |   |-- api/                 # FastAPI routers
|   |   `-- tests/               # pytest unit tests
|   |-- requirements.txt
|   `-- .env.example
|-- frontend/
|   |-- src/
|   |   |-- app/                 # Next.js App Router pages
|   |   |-- components/          # Per-module React components + charts
|   |   |-- lib/                 # TypeScript types, API client
|   |   `-- hooks/               # useAuditSession context
|   `-- package.json
`-- README.md
\end{lstlisting}

\subsection{Request Flow}

Each module follows the same pattern: the frontend collects structured input
through a form component, POSTs it to the corresponding FastAPI endpoint, the
backend validates it against a Pydantic schema, dispatches to a pure
calculation function in \texttt{app/core/}, and returns a structured JSON
response that the frontend renders as tables and charts (Recharts). The
optional AI-insights endpoint and the PDF report endpoint both consume the
same calculation result objects, ensuring the report and the on-screen
results are always consistent.

\subsection{Configuration}

Backend configuration is loaded from a \texttt{.env} file via
\texttt{pydantic-settings} (\texttt{app/config.py}). Key settings include:

\begin{table}[h]
\centering
\small
\begin{tabular}{@{}p{6.2cm}p{2.6cm}p{4.5cm}@{}}
\toprule
\textbf{Setting} & \textbf{Default} & \textbf{Purpose} \\
\midrule
\code{cors\_origins} & localhost:3000 & Allowed frontend origin(s) \\
\code{default\_tariff\_per\_kwh} & 0.12 & Default electricity tariff \\
\code{default\_co2\_factor\_kg\_per\_kwh} & 0.233 & Grid CO\textsubscript{2} intensity \\
\code{default\_lifetime\_years} & 20 & Motor lifetime for savings analysis \\
\code{default\_vfd\_efficiency} & 0.97 & Default VFD drive efficiency \\
\code{default\_vfd\_headroom} & 1.10 & VFD sizing headroom factor \\
\code{pq\_thd\_voltage\_limit\_pct} & 5.0 & Voltage THD status-flag limit \\
\code{pq\_thd\_current\_limit\_pct} & 8.0 & Current THD status-flag limit \\
\code{pq\_power\_factor\_min} & 0.95 & Power-factor status-flag limit \\
\code{pq\_current\_unbalance\_pct} & 3.0 & Current-unbalance status-flag limit \\
\code{gemini\_api\_key} & (empty) & Enables AI-insights endpoints \\
\bottomrule
\end{tabular}
\caption{Selected backend configuration settings.}
\end{table}

% ══════════════════════════════════════════════════════════════════════════
\section{Installation and Quick Start}
% ══════════════════════════════════════════════════════════════════════════

\subsection{Backend}

\begin{lstlisting}[style=code, language=bash]
cd motoraudit/backend
python3 -m venv .venv && source .venv/bin/activate   # Windows: .venv\Scripts\activate
pip install -r requirements.txt
cp .env.example .env
uvicorn app.main:app --reload --port 8000
# API docs -> http://localhost:8000/docs
\end{lstlisting}

\subsection{Frontend}

\begin{lstlisting}[style=code, language=bash]
cd motoraudit/frontend
npm install
cp .env.example .env.local
npm run dev
# App -> http://localhost:3000
\end{lstlisting}

\subsection{Tests}

\begin{lstlisting}[style=code, language=bash]
cd motoraudit/backend
pytest -v
\end{lstlisting}

Unit tests cover the DSP engine (\texttt{test\_dsp.py}), the replacement
calculator (\texttt{test\_replacement.py}), and the VFD calculator
(\texttt{test\_vfd.py}).

% ══════════════════════════════════════════════════════════════════════════
\section{Module 1: Motor Replacement \& Payback}
\label{sec:module1}
% ══════════════════════════════════════════════════════════════════════════

\subsection{Engineering Rationale}

A motor that runs 6\,000 h/yr at 75\% load for 20 years consumes roughly ten
times its own purchase price in electricity. Moving from an IE1 to an IE3
efficiency class typically saves 2--4\% of input energy at the operating
load point; on large machines running long hours, this translates into
substantial annual savings.

\subsection{Motor Efficiency Model}

The IEC 60034-30-1:2014 standard (Table 1) tabulates minimum guaranteed
efficiency $\eta$ at four standard load points: 25\%, 50\%, 75\%, and 100\%
of rated load, for each IE class (IE1--IE4), motor pole count, and rated
power. The efficiency at an arbitrary load factor $\mathrm{LF}$ is obtained
by fitting a least-squares quadratic polynomial through the four tabulated
$(\mathrm{LF}, \eta)$ points and evaluating it at the requested load factor:

\begin{equation}
\eta(\mathrm{LF}) = a\,\mathrm{LF}^2 + b\,\mathrm{LF} + c
\end{equation}

where $a$, $b$, $c$ are obtained from a quadratic least-squares fit through
the four IEC points $\mathrm{LF} \in \{0.25, 0.50, 0.75, 1.00\}$. A quadratic
is adequate because the $\eta$-vs-$\mathrm{LF}$ curve of an induction motor
is smooth and unimodal over this range. The result is clipped to $[0.30,
1.00]$ as a numerical safeguard.

If the user instead supplies a single full-load efficiency value
(\code{efficiency\_override\_pct}) rather than an IE class, a parametric
Gaussian-peak model is used instead of the IEC table:

\begin{equation}
\eta(\mathrm{LF}) = \eta_{\mathrm{peak}} \cdot
    \exp\!\big(-k\,(\mathrm{LF} - \mathrm{LF}_{\mathrm{peak}})^2\big)
\end{equation}

with $\mathrm{LF}_{\mathrm{peak}} = 0.75$ and $k$ chosen so that
$\eta(1.0) \approx 0.985\,\eta_{\mathrm{peak}}$, i.e.\
$k = -\ln(0.985)/0.0625 \approx 0.243$. The requested load factor is clamped
to $[0.05, 1.05]$ before evaluation in both models.

\textbf{Data source.} IEC 60034-30-1:2014 Table 1 minimum efficiency values,
4-pole 50\,Hz motors, IE1--IE4, 4--110\,kW. These are minimum \emph{guaranteed}
values --- name-plate efficiencies of specific commercial models are
typically 1--2\% higher.

\subsection{Energy and Cost Calculations}

Electrical input power at the operating load factor is the shaft output
power divided by efficiency:

\begin{equation}
P_{\mathrm{shaft}} = P_{\mathrm{rated}} \cdot \mathrm{LF}, \qquad
P_{\mathrm{in}} = \frac{P_{\mathrm{shaft}}}{\eta(\mathrm{LF})}
\end{equation}

Annual energy consumption:

\begin{equation}
E_{\mathrm{annual}} = \frac{P_{\mathrm{rated}} \cdot \mathrm{LF}}{\eta(\mathrm{LF})} \cdot H
\qquad \text{[kWh/yr]}
\end{equation}

where $H$ is the annual operating hours. The same load factor and rated
power are used for both the existing and replacement motor, so that only the
efficiency term changes between the two scenarios.

Annual savings from replacing the existing motor ($E_{\mathrm{old}}$) with the
new one ($E_{\mathrm{new}}$):

\begin{align}
\Delta E &= E_{\mathrm{old}} - E_{\mathrm{new}} && \text{[kWh/yr]} \\
\Delta \mathrm{Cost} &= \Delta E \cdot \mathrm{tariff} && \text{[currency/yr]} \\
\Delta \mathrm{CO_2} &= \Delta E \cdot \mathrm{co2\_factor} && \text{[kg CO\textsubscript{2}/yr]}
\end{align}

Simple (undiscounted) payback period:

\begin{equation}
\mathrm{Payback} = \frac{C_{\mathrm{motor}}}{\Delta \mathrm{Cost}} \qquad \text{[years]}
\end{equation}

Lifetime net savings over the assumed motor lifetime $L$ (also undiscounted,
i.e.\ no NPV/IRR):

\begin{equation}
\mathrm{LNS} = \Delta \mathrm{Cost} \cdot L - C_{\mathrm{motor}}
\end{equation}

\subsection{API Contract Summary}

\textbf{Request} (\code{ReplacementRequest}): \code{existing\_motor} and
\code{replacement\_motor} objects (each specifying either \code{ie\_class} or
\code{efficiency\_override\_pct}, plus \code{rated\_power\_kw} and
\code{load\_factor} for the existing motor, and \code{cost} for the
replacement), \code{operating\_hours\_per\_year}, \code{tariff\_per\_kwh},
optional \code{co2\_factor\_kg\_per\_kwh}, and \code{lifetime\_years}.

\textbf{Response} (\code{ReplacementResponse}): per-motor operating points
(\code{existing}, \code{replacement}), a \code{savings} summary, an
\code{efficiency\_curve} for charting $\eta$ vs.\ load factor for both
motors, a \code{payback\_timeline} (cumulative savings and net position by
year), and a list of human-readable \code{assumptions} describing the exact
inputs used.

% ══════════════════════════════════════════════════════════════════════════
\section{Module 2: VFD Sizing \& Savings}
\label{sec:module2}
% ══════════════════════════════════════════════════════════════════════════

\subsection{Engineering Rationale}

Pump and fan systems are among the largest categories of industrial
electricity use. When flow is controlled by throttling a valve or damping a
fan, the motor continues to run near full speed and the excess energy is
wasted across the restriction. A variable frequency drive (VFD) instead
reduces motor speed to match actual flow demand. Because shaft power on
centrifugal loads follows the \textbf{affinity law} (power $\propto$
speed\textsuperscript{3}), running at 70\% speed consumes only $0.7^3
\approx 34\%$ of full-speed power.

\subsection{Load Models}

\textbf{Centrifugal loads (pump / fan)} obey the cubic affinity law:

\begin{equation}
P_{\mathrm{shaft}}(n) = P_{\mathrm{rated}} \cdot \left(\frac{n}{n_{\mathrm{rated}}}\right)^3
\end{equation}

\textbf{Constant-torque loads (conveyor)} scale linearly with speed:

\begin{equation}
P_{\mathrm{shaft}}(n) = P_{\mathrm{rated}} \cdot \frac{n}{n_{\mathrm{rated}}}
\end{equation}

In both cases, the electrical input through the VFD-driven motor is:

\begin{equation}
P_{\mathrm{elec}} = \frac{P_{\mathrm{shaft}}}{\eta_{\mathrm{motor}} \cdot \eta_{\mathrm{VFD}}}
\end{equation}

\textbf{Throttling baseline.} In the absence of a VFD, the motor is modelled
as drawing a constant electrical input of $P_{\mathrm{rated}} /
\eta_{\mathrm{motor}}$ at every operating point, regardless of actual flow
demand. This is a deliberately conservative assumption (some motors draw
marginally less at high slip) and is the standard convention for VFD savings
calculations.

\subsection{Annual Energy over a Load Profile}

The user supplies a load profile as a list of $(n_i, t_i)$ pairs, where $n_i$
is the speed fraction and $t_i$ is the fraction of annual operating hours
spent at that speed ($\sum_i t_i = 1$, enforced within $\pm 0.02$ tolerance).
For each profile point $i$:

\begin{equation}
E_i = P_i \cdot t_i \cdot H \qquad \text{[kWh/yr]}
\end{equation}

Total baseline and VFD energy are the sums of $E_i$ over all profile points
using the respective $P_i$ model above. Savings and payback follow the same
form as Module 1:

\begin{equation}
\Delta E = E_{\mathrm{baseline}} - E_{\mathrm{VFD}}, \qquad
\mathrm{Payback} = \frac{C_{\mathrm{VFD}}}{\Delta E \cdot \mathrm{tariff}}
\end{equation}

\subsection{VFD Sizing}

The recommended VFD power rating includes a headroom margin to cover
starting-torque boost and thermal derating:

\begin{equation}
P_{\mathrm{VFD}} = P_{\mathrm{motor}} \cdot \mathrm{headroom\_factor}
\end{equation}

with a default headroom factor of 1.10 (10\%).

\subsection{API Contract Summary}

\textbf{Request} (\code{VFDRequest}): \code{motor\_power\_kw},
\code{load\_type} (\code{pump} / \code{fan} / \code{conveyor}),
\code{motor\_efficiency}, \code{vfd\_efficiency}, a \code{load\_profile}
list of $\{$\code{speed\_fraction}, \code{time\_fraction}$\}$ points (1--20
entries, time fractions summing to 1.0), \code{operating\_hours\_per\_year},
\code{tariff\_per\_kwh}, \code{vfd\_cost}, and \code{headroom\_factor}.

\textbf{Response} (\code{VFDResponse}): \code{recommended\_vfd\_kw}, a
per-point \code{profile\_analysis} breakdown, aggregate \code{totals}
(baseline/VFD annual energy, savings, payback), a \code{power\_vs\_speed}
curve (10 points from 10\% to 100\% speed) for charting, and a list of
\code{assumptions}.

% ══════════════════════════════════════════════════════════════════════════
\section{Module 3: Power Quality Analyzer}
\label{sec:module3}
% ══════════════════════════════════════════════════════════════════════════

\subsection{Engineering Rationale}

Non-linear loads --- VFDs, rectifiers, UPS systems --- inject harmonic
currents that increase I\textsuperscript{2}R losses, overheat motors and
transformers, and can cause nuisance protective-relay tripping. Voltage
unbalance above roughly 2\% reduces motor efficiency and shortens motor
life. Low power factor increases apparent current draw and, in many
tariff structures, incurs utility penalty charges. Module 3 analyzes
voltage and current waveforms (three-phase or single-phase) to quantify
these effects.

\subsection{Signal Processing Pipeline}

Each channel (\code{ia}, \code{ib}, \code{ic}, \code{va}, \code{vb},
\code{vc}) is processed with a Hanning-windowed FFT to reduce spectral
leakage:

\begin{equation}
X[k] = \mathrm{FFT}\big(x[n] \cdot w[n]\big), \qquad
w[n] = 0.5\left(1 - \cos\!\frac{2\pi n}{N-1}\right)
\end{equation}

Peak amplitude at each frequency bin is corrected for the window's energy
loss, and converted to RMS:

\begin{equation}
A_{\mathrm{peak}}[k] = \frac{2\,|X[k]|}{\sum_n w[n]}, \qquad
A_{\mathrm{rms}}[k] = \frac{A_{\mathrm{peak}}[k]}{\sqrt{2}}
\end{equation}

For each harmonic order $h = 1, 2, \ldots$, the analyzer searches a $\pm
3$\,Hz window around the target frequency $h \cdot f_0$ (to tolerate slight
frequency drift between the signal and the sample-rate grid) and takes the
maximum RMS amplitude found as $I_h$ (or $V_h$).

\subsection{Total Harmonic Distortion (THD)}

Following the IEEE 519-2022 definition:

\begin{equation}
\mathrm{THD} = \frac{\sqrt{\displaystyle\sum_{h=2}^{N} I_h^2}}{I_1} \times 100\%
\end{equation}

where $I_1$ is the RMS amplitude of the fundamental and $I_h$ are the RMS
amplitudes of the harmonic components (orders 1--15 are always retained;
higher orders are retained only if their amplitude exceeds 0.3\% of the
fundamental).

\subsection{RMS and Peak}

\begin{equation}
X_{\mathrm{rms}} = \sqrt{\frac{1}{N}\sum_{i=1}^{N} x_i^2}, \qquad
X_{\mathrm{peak}} = \max_i |x_i|
\end{equation}

\subsection{Power Metrics}

Given synchronized voltage $v(t)$ and current $i(t)$ samples for one phase:

\begin{align}
P &= \overline{v \cdot i} && \text{(active power)} \\
S &= V_{\mathrm{rms}} \cdot I_{\mathrm{rms}} && \text{(apparent power)} \\
Q &= \sqrt{\max(S^2 - P^2,\ 0)} && \text{(reactive power)} \\
\mathrm{PF} &= \frac{P}{S} && \text{(true power factor)}
\end{align}

The \textbf{displacement power factor} $\cos\varphi$ is computed from the
phase angle between the fundamental voltage and current phasors (obtained
from the FFT bin nearest the detected fundamental frequency):

\begin{equation}
\varphi = \angle V_1 - \angle I_1, \qquad \cos\varphi = |\cos(\varphi)|
\end{equation}

Note that $\mathrm{PF} < \cos\varphi$ whenever harmonic distortion is
present, since $\mathrm{PF}$ (true power factor) includes the effect of
distortion power, while $\cos\varphi$ reflects only the fundamental-frequency
displacement.

\textbf{Reactive power correction.} When $\mathrm{PF}$ falls below the
configured target (default 0.95), the required shunt capacitor bank rating
to reach the target is:

\begin{equation}
Q_{\mathrm{corr}} = P \cdot \big(\tan\varphi_{\mathrm{actual}} - \tan\varphi_{\mathrm{target}}\big)
\qquad \text{[kVAr]}
\end{equation}

\subsection{Phase Unbalance (NEMA MG-1)}

For a three-phase system with RMS phase currents $I_a, I_b, I_c$:

\begin{equation}
\bar{I} = \frac{I_a + I_b + I_c}{3}, \qquad
\mathrm{UB\%} = \frac{\max_x |I_x - \bar{I}|}{\bar{I}} \times 100
\end{equation}

The same formula is applied to phase voltages when all three voltage
channels are present.

\subsection{Status Flags and Recommendations}

Each computed metric is compared against a configurable threshold
(\code{PQThresholds}: \code{thd\_voltage\_pct}, \code{thd\_current\_pct},
\code{power\_factor\_min}, \code{current\_unbalance\_pct}) and classified as
\textbf{PASS}, \textbf{WARN} (within 20\% over the limit, or within 2\% under
a ``low-is-bad'' limit such as power factor), or \textbf{FAIL}. A
rule-based recommendation engine then produces plain-language corrective
guidance --- e.g.\ identifying dominant 5th/7th harmonics as indicative of a
6-pulse VFD or rectifier load and suggesting a passive or active harmonic
filter, or sizing a capacitor bank from the computed $Q_{\mathrm{corr}}$.

These thresholds and recommendations are \textbf{advisory} engineering
estimates approximating IEEE 519-2022 general system limits; they do not
constitute legal or regulatory compliance certification.

\subsection{Preset Signals}

For demonstration and testing without uploaded data, four synthetic
waveform presets are available (sampled at 10\,kHz, 50\,Hz fundamental,
0.2\,s duration):

\begin{table}[h]
\centering
\begin{tabular}{@{}ll@{}}
\toprule
\textbf{Preset} & \textbf{Description} \\
\midrule
\code{clean\_sinusoid}     & Pure 50\,Hz, 42\,A / 230\,V --- zero-THD reference \\
\code{vfd\_harmonics}      & 6-pulse VFD current: 5th (14\%), 7th (9\%), 11th (5\%), 13th (3\%) \\
\code{unbalanced\_3phase}  & Three-phase with phase B $-8\%$, phase C $+6\%$ unbalance \\
\code{low\_power\_factor}  & Single-phase, $37^\circ$ lag $\rightarrow$ PF $\approx 0.80$ \\
\bottomrule
\end{tabular}
\caption{Built-in Module 3 waveform presets.}
\end{table}

\subsection{API Contract Summary}

\textbf{Request} (\code{PQRequest}): \code{source} (\code{json} or
\code{preset}), optional \code{preset\_name}, \code{sample\_rate\_hz},
\code{fundamental\_freq\_hz}, \code{thresholds}, and optional \code{data}
(a \code{WaveformData} object with \code{ia} required and \code{ib}, \code{ic},
\code{va}, \code{vb}, \code{vc} optional, each a list of $\geq 64$ samples of
equal length). A dedicated upload endpoint additionally accepts the waveform
as a CSV file with columns \code{time, ia[, ib, ic, va, vb, vc]}.

\textbf{Response} (\code{PQResponse}): \code{duration\_s},
\code{actual\_fundamental\_hz} (detected from the \code{ia} or \code{va}
FFT, not assumed), per-channel \code{channels} analysis (RMS, peak, THD,
harmonic table), optional \code{power\_metrics}, optional
\code{phase\_unbalance}, a \code{status\_table} of \code{StatusFlag} entries,
a list of \code{recommendations}, and downsampled \code{waveform\_preview}
and \code{fft\_preview} data (log-magnitude in dBFS, 0--2\,kHz) for
charting.

% ══════════════════════════════════════════════════════════════════════════
\section{AI-Generated Engineering Insights}
% ══════════════════════════════════════════════════════════════════════════

Each module has a corresponding \texttt{/api/ai/*} endpoint that sends the
structured calculation result to Google's Gemini model
(\texttt{gemini-3-flash-preview}) with a role-primed prompt (``a senior
energy engineer with 20 years of experience'') and a fixed three-paragraph
output template tailored to that module:

\begin{itemize}
    \item \textbf{Replacement:} business-case framing (is the payback
          attractive?), sensitivity analysis (which input has the most
          leverage?), and pre-commissioning checks.
    \item \textbf{VFD:} explanation of why the cubic affinity-law savings
          are physically expected, commonly overlooked installation risks
          (PWM dV/dt insulation stress, minimum-speed cooling limits,
          harmonic injection, bearing currents), and further optimization
          strategies.
    \item \textbf{Power Quality:} likely equipment source of the observed
          harmonic signature, quantified consequences for connected
          equipment and billing, and a prioritized corrective sequence.
    \end{itemize}

The Gemini client is lazily initialized; if \code{GEMINI\_API\_KEY} is not
configured, the AI-insights endpoints return \code{503 Service Unavailable}
while the rest of the application continues to function normally.

% ══════════════════════════════════════════════════════════════════════════
\section{PDF Report Generation}
% ══════════════════════════════════════════════════════════════════════════

The \apiendpoint{POST}{/api/report/summary} endpoint accepts the results of any
combination of the three modules and returns a consolidated JSON summary.
The \apiendpoint{POST}{/api/report/pdf} endpoint performs the same
consolidation and streams a formatted PDF audit report, generated with
ReportLab, suitable for direct delivery to a client or inclusion in a plant
energy-audit file.

% ══════════════════════════════════════════════════════════════════════════
\section{REST API Reference}
% ══════════════════════════════════════════════════════════════════════════

All endpoints are served under the \code{/api} prefix. Interactive
Swagger UI documentation is available at
\href{http://localhost:8000/docs}{\code{http://localhost:8000/docs}}
(ReDoc at \code{/redoc}) whenever the backend is running.

\begin{longtable}{@{}p{1.7cm}p{5.3cm}p{7.3cm}@{}}
\toprule
\textbf{Method} & \textbf{Path} & \textbf{Description} \\
\midrule
\endhead
GET  & \code{/api/health} & Service health check \\
GET  & \code{/api/motors} & List seeded IEC motors \\
GET  & \code{/api/replacement/demo} & Demo preset inputs for Module 1 \\
POST & \code{/api/replacement/calculate} & Run Module 1 calculation \\
GET  & \code{/api/vfd/demo} & Demo preset inputs for Module 2 \\
POST & \code{/api/vfd/calculate} & Run Module 2 calculation \\
GET  & \code{/api/power-quality/presets} & List available Module 3 preset names \\
POST & \code{/api/power-quality/analyze} & Run Module 3 analysis on JSON waveform data \\
POST & \code{/api/power-quality/upload} & Run Module 3 analysis on an uploaded CSV file \\
POST & \code{/api/report/summary} & Consolidated JSON summary across modules \\
POST & \code{/api/report/pdf} & Consolidated PDF report download \\
POST & \code{/api/ai/replacement} & AI insights for a Module 1 result \\
POST & \code{/api/ai/vfd} & AI insights for a Module 2 result \\
POST & \code{/api/ai/power-quality} & AI insights for a Module 3 result \\
\bottomrule
\caption{MotorAudit REST API endpoints.}
\end{longtable}

% ══════════════════════════════════════════════════════════════════════════
\section{Frontend Application}
% ══════════════════════════════════════════════════════════════════════════

The frontend is a Next.js (App Router) application with one route per
module: \code{/replacement}, \code{/vfd}, \code{/power-quality}, plus a
\code{/simulation} route hosting an interactive motor simulation
(\code{MotorSimulation.tsx}). Each module route pairs an input form
component (e.g.\ \code{ReplacementForm.tsx}, \code{VFDForm.tsx},
\code{PQForm.tsx}) with a results component (e.g.\
\code{ReplacementResults.tsx}) and one or more Recharts-based visualizations
(e.g.\ \code{EfficiencyVsLoadChart.tsx}, \code{PaybackTimelineChart.tsx},
\code{PowerVsSpeedChart.tsx}, \code{FFTSpectrumChart.tsx},
\code{WaveformChart.tsx}). Shared components include \code{FormulaBox.tsx}
(renders the governing formula alongside results), \code{AIInsights.tsx}
(displays the Gemini-generated commentary), and
\code{AuditSummaryPanel.tsx}. Cross-module state --- allowing a user to run
all three modules and then generate one consolidated report --- is held in
the \code{useAuditSession} React context (\code{hooks/useAuditSession.ts}),
and all HTTP calls to the backend are centralized in
\code{lib/api.ts} with request/response types declared in \code{lib/types.ts}.

% ══════════════════════════════════════════════════════════════════════════
\section{Example Results (Demo Presets)}
% ══════════════════════════════════════════════════════════════════════════

\subsection{Module 1 --- 37\,kW IE1 $\to$ IE3, 75\% Load, 6\,000 h/yr, \$0.12/kWh}

\begin{itemize}[nosep]
    \item Annual energy saved: \textbf{$\approx$ 8\,280 kWh}
    \item Annual cost saved: \textbf{$\approx$ \$994}
    \item Simple payback: \textbf{$\approx$ 4.2 years} (at \$4\,200 motor cost)
    \item CO\textsubscript{2} reduction: \textbf{$\approx$ 1\,929 kg/yr} (at 0.233 kg/kWh)
\end{itemize}

\subsection{Module 2 --- 55\,kW Pump VFD, Profile 10/30/40/20\% at 100/85/70/50\% Speed, 7\,000 h/yr}

\begin{itemize}[nosep]
    \item Baseline: \textbf{$\approx$ 415\,800 kWh/yr} (throttling)
    \item VFD: \textbf{$\approx$ 246\,000 kWh/yr}
    \item Saved: \textbf{$\approx$ 170\,000 kWh/yr} ($\approx$ 41\% reduction)
    \item Payback: \textbf{$\approx$ 0.4 years} at \$7\,500 VFD cost
\end{itemize}

\subsection{Module 3 --- VFD Harmonics Preset}

\begin{itemize}[nosep]
    \item Current THD: \textbf{$\approx$ 18\%} $\to$ \textbf{FAIL} ($>$ 8\% limit)
    \item Voltage THD: \textbf{$\approx$ 1\%} $\to$ \textbf{PASS}
    \item Power factor: \textbf{0.999} $\to$ \textbf{PASS} (clean voltage, no reactive component in demo)
    \item Recommendation: passive harmonic filter targeting the 5th and 7th harmonics
\end{itemize}

% ══════════════════════════════════════════════════════════════════════════
\section{Assumptions and Limitations}
\label{sec:assumptions}
% ══════════════════════════════════════════════════════════════════════════

\begin{itemize}
    \item Motor efficiency data is sourced from IEC 60034-30-1:2014 minimum
          guaranteed values (4-pole, 50\,Hz). Real motors are typically
          1--2\% more efficient than the tabulated minimum.
    \item Efficiency interpolation uses a quadratic polynomial fit through
          four IEC load points; accuracy degrades for load factors outside
          the 25--100\% range.
    \item VFD savings assume a quadratic (resistance-dominated) system
          curve. A system with a significant static head component will
          show lower actual savings than predicted.
    \item Power-quality thresholds are advisory approximations of IEEE
          519-2022 general limits. No legal or regulatory compliance
          certification is implied by PASS/WARN/FAIL status flags.
    \item All financial calculations are undiscounted --- no net present
          value (NPV) or internal rate of return (IRR) is computed. Users
          requiring rigorous investment analysis should apply an
          appropriate discount rate to the reported cash flows.
    \item Results are labelled as \textbf{engineering estimates} throughout
          the user interface and generated reports.
\end{itemize}

% ══════════════════════════════════════════════════════════════════════════
\section{References}
% ══════════════════════════════════════════════════════════════════════════

\begin{enumerate}
    \item IEC 60034-30-1:2014, \textit{Rotating electrical machines --- Part
          30-1: Efficiency classes of line operated AC motors (IE code)}.
    \item IEEE Std 519-2022, \textit{IEEE Standard for Harmonic Control in
          Electric Power Systems}.
    \item NEMA MG-1, \textit{Motors and Generators}, National Electrical
          Manufacturers Association.
\end{enumerate}

\end{document}
